\section{Service Agreements and Their Laminar Allocation Form}\label{sec:model}
\subsection{Public primitives and accepted continuations}
Let $G=(\mathcal V,\mathcal E)$ be a finite, rooted, acyclic public graph. Vertices include date; all are reachable. Write $N=|\mathcal V|$, $M=|\mathcal E|$, and let $p_{vw}>0$ be the exogenous probability on edge $(v,w)$, with probabilities summing to one at a nonterminal vertex. The discount factor is $0<\beta\leq1$. An occurrence $h$ is a root-to-vertex history in the finite unfolding $\mathcal T$; its public vertex is $v(h)$ and its discounted reach weight is $w_h=\beta^{t(h)}\Pr(h)>0$. Write $j\succeq h$ for a descendant, including $h$, and $H=|\mathcal T|$.

At occurrence $h$ of $v$, the provider chooses a tier $x_h\in[l_v,h_v]\subseteq[0,1]$. The payment coefficient $a_v>0$ is fixed, and the operating reward is $g_v(x)=r_vx-q_vx^2/2$, with $q_v>0$. Customer utility is a publicly specified service-benefit term minus $a_vx$. For example, fixed stock $s_z$ yields the stage utility
\begin{equation}\label{eq:customer-r30}
 u(z,s_z,x)=\nu\mathbb E\min(s_z,D_z)-a_zx,
 \qquad a_z=r_z^{\rm premium}-\kappa\mathbb E(D_z-s_z)^+>0.
\end{equation}
Stock may be stipulated, or identified by the service and cost commitments in the retained stock-identification results. The customer can leave before each modeled review; the provider commits to its contingent protocol. Retaining the customer requires the conditional service benefit minus remaining payment to exceed the public outside continuation. Assume their difference is a vertex-sufficient cap $B_v$. The accepted optimization is
\begin{align}
 J(b_0)=\max_x\quad &\sum_{h\in\mathcal T}w_hg_{v(h)}(x_h) \label{eq:primal-r30}\\
 \text{subject to}\quad
 &Y_h:=\sum_{j\succeq h}\frac{w_j}{w_h}a_{v(j)}x_j\leq B_{v(h)}\quad(h\in\mathcal T),\qquad Y_0=b_0,\nonumber\\
 &l_{v(h)}\leq x_h\leq h_{v(h)}.\nonumber
\end{align}
The root has a cap as well as the equality promise. No privately informed reports, endogenous termination choice, or policy-dependent transition probabilities are part of this institution. The no-switching restriction is essential to the scalar-price result; the retained general model continues to allow switching and is not replaced by \eqref{eq:primal-r30}.

The feasible remaining-payment interval is computed without unfolding:
\begin{equation}\label{eq:intervals-r30}
 L_v=a_vl_v+\beta\sum_wp_{vw}L_w,\qquad
 U_v=\min\left\{B_v,a_vh_v+\beta\sum_wp_{vw}U_w\right\}.
\end{equation}
A subtree is feasible precisely when its successors are feasible and $L_v\leq U_v$. Throughout, $b_0\in[L_0,U_0]\subseteq(-\infty,B_0]$. Convexity fills every point between the two extreme payments. Compactness gives an optimum and strict concavity gives a unique occurrence-indexed action vector.

\subsection{An exact mapping, not an analogy}
\begin{proposition}[Laminar allocation equivalence]\label{prop:laminar}
Define $z_h=w_ha_{v(h)}x_h$. Problem~\eqref{eq:primal-r30} is exactly the separable concave allocation problem
\begin{align}
 \max_z\quad &\sum_h\left[\frac{r_{v(h)}}{a_{v(h)}}z_h-
 \frac{q_{v(h)}}{2w_ha_{v(h)}^2}z_h^2\right],\label{eq:laminar}\\
 \sum_{j\succeq h}z_j&\leq w_hB_{v(h)},\qquad \sum_jz_j=b_0,\nonumber\\
 w_ha_{v(h)}l_{v(h)}&\leq z_h\leq w_ha_{v(h)}h_{v(h)}.\nonumber
\end{align}
The constrained descendant sets are laminar. The change of variables is invertible and preserves objective values and optimal policies.
\end{proposition}
\begin{proof}
Substitution gives each term and row of \eqref{eq:laminar}. Two rooted descendant sets in a tree are either disjoint or one contains the other. Positivity of $w_ha_{v(h)}$ makes the coordinate map bijective and preserves every inequality direction.\end{proof}

Thus the expanded primal lies in the tree-constrained allocation lineage of \citet{Mjelde1983} and \citet{Tang1990}. Tang reports at most twice the number of variables in one-variable convex subproblems. That count is not a count of elementary rational operations for a complete parametric response family. The chain-nested allocation results of \citet{Vidal2016,SchootUiterkamp2021,Wu2021} concern related but different input structures. In particular, a complexity bound for an ascending chain cannot simply be assigned to arbitrary branching descendant sets.

After subtracting lower allocations, the packing region is a laminar polymatroid, and the exact total selects a base of its truncation (Electronic Companion, Section~\ref{ec:polymatroid}). This connects the model to the broader allocation literature, including \citet{FedergruenGroenevelt1986} on discrete greedy allocation and \citet{Groenevelt1991} on separable concave maximization over polymatroids. These connections identify inherited feasible-set geometry; they do not by themselves provide a compact-DAG response-size or writable-memory bound. Nor do the discrete polynomial bounds stated in that literature's abstracts establish the binary-input bound proved here for the continuous quadratic response construction.

\paragraph{What memoization does and does not change.}
An occurrence-indexed marginal reduction may represent its subtree payment in the unnormalized units $z$. Its entire response is the occurrence weight times a common conditional response, proved below. Dividing by that weight makes occurrences of the same public vertex identical. Memoizing those normalized responses therefore gives our graph recursion. There is no additional stochastic obstruction or different first-order condition. The quantitative questions are whether complete responses have a bounded common knot universe, whether exact coefficients have polynomial bit length, and how the resulting policy can be minimized as an execution machine. Theorems~\ref{thm:quotient}--\ref{thm:bit} and \ref{thm:minimal} answer those questions directly. We do not assert that scalar marginal allocation, subtree reduction, or the principle of memoization originated here.

An explicit-tree implementation has input size $H$ and must at least read its $H$ local records; a compact representation has $N$ vertices and $M$ edges. This is an input/output comparison, not a lower bound against every algorithm that accepts a compact graph. The independent tree baseline below implements a specified exact reduction; it is not presented as the original authors' implementation or as the fastest available laminar algorithm.

\subsection{Relation to contractual states and policy graphs}
Promised continuation utility is a standard recursive state in dynamic contracting \citep{SpearSrivastava1987,ThomasWorrall1988,ChenSunXiao2020}. Our remaining payment is a primal contractual state; an incoming payment price is its dual counterpart. The public-information institution does not solve the incentive or adverse-event allocation problems in \citet{ChenSunXiao2020,LiangSunTangZhang2023}. It instead determines what accepted service adaptation requires after public continuation constraints have been imposed.

Policy graphs and stochastic decomposition already avoid naive scenario enumeration \citep{PereiraPinto1991,Girardeau2015,Dowson2020}. We compare an explicit primal promise-grid policy graph below, rather than treat the mere use of a public graph as a contribution. Our exact closure depends on one scalar additive commitment and separable quadratics; general convex policy-graph decomposition has a broader domain. The supporting cuts for fixed tables are also instances of standard dual decomposition \citep{Benders1962,BoydVandenberghe2004}, not a new general finite-convergence theorem.
